% GENERATED FILE. Edit canonical lessons, then run npm run generate:books.
% canonical-source-hash: fnv1a64:ba6db3c3628d2c92
% canonical-lessons: SA-C534-madhyasthah, SA-C534-samvid, SA-C534-angikaroti, SA-C534-pratyakhyati, SA-C534-sammatih

\chapter{Coming to Terms}
\label{ch:sa-coming-to-terms}
\hlchaptermodality{534}

\begin{chapteropening}
\textbf{By the end of this chapter:} \emph{I can talk about agreements, go-betweens and saying yes or no.}

Complete the last lesson of chapter 534: I can talk about agreements, go-betweens and saying yes or no.
\end{chapteropening}

\section[\texorpdfstring{\sk{मध्यस्थः} (madhyasthaḥ)}{मध्यस्थः (madhyasthaḥ)}]{\sk{मध्यस्थः} (madhyasthaḥ) — a go-between}

\label{lesson:SA-C534-madhyasthah}



\begin{quote}
Before the new one: say the Sanskrit for moves out of the way, then the Sanskrit for trouble.
\end{quote}

\subsection*{You'll want to know: \sk{मध्यस्थः}}
\textbf{\sk{मध्यस्थः}} — \emph{madhyasthaḥ} — \textquotedblleft{}a go-between\textquotedblright{}.

\begin{quote}
\sk{यदि} \sk{विवादः} \sk{अस्ति}, \sk{तर्हि} \sk{मध्यस्थः} \sk{समाधानं} \sk{ददातु।}
\end{quote}

\begin{quote}
\emph{yadi vivādaḥ asti, tarhi madhyasthaḥ samādhānaṃ dadātu.}
\end{quote}

\begin{quote}
If there is a dispute, let a go-between offer a way out.
\end{quote}

A masculine noun.

\begin{cousinweb}
\textbf{\sk{मध्यस्थः}} (\emph{madhyasthaḥ}) is built from \emph{madhya}, \textquotedblleft{}the middle\textquotedblright{}, and \emph{stha}, \textquotedblleft{}standing\textquotedblright{}, from the root √sthā; MW gives \emph{madhyastha} as standing between two parties, a mediator, and as impartial (MW).
\end{cousinweb}

\begin{grammarlens}[title={What you've built}]
One more word for saying what happens if.
\end{grammarlens}

\subsection*{Your turn}
\begin{itemize}
\raggedright
  \item \emph{Say it:} \emph{madhyasthaḥ}
  \item \emph{Say it:} \emph{madhyasthaḥ}, once more
  \item \emph{Say it:} \emph{yadi vivādaḥ asti, tarhi madhyasthaḥ samādhānaṃ dadātu.} — If there is a dispute, let a go-between offer a way out.
  \item \emph{Recall:} say the Sanskrit for moves out of the way, then the Sanskrit for trouble, then say \emph{madhyasthaḥ} again
\end{itemize}

\begin{tcolorbox}[breakable,colback=teal!4,colframe=teal!35!black,title={Before you move on}]
What does \sk{मध्यस्थः} mean? (\textquotedblleft{}A go-between\textquotedblright{}.) Say it once more.
\end{tcolorbox}
\section[\texorpdfstring{\sk{संविद्} (saṃvid)}{संविद् (saṃvid)}]{\sk{संविद्} (saṃvid) — an agreement}

\label{lesson:SA-C534-samvid}



\begin{quote}
Before the new one: say the Sanskrit for trouble, then the Sanskrit for a go-between.
\end{quote}

\subsection*{You'll want to know: \sk{संविद्}}
\textbf{\sk{संविद्}} — \emph{saṃvid} — \textquotedblleft{}an agreement\textquotedblright{}.

\begin{quote}
\sk{यदि} \sk{भवान्} \sk{तेन} \sk{सह} \sk{संविदं} \sk{करोति}, \sk{तर्हि} \sk{कार्यं} \sk{भविष्यति।}
\end{quote}

\begin{quote}
\emph{yadi bhavān tena saha saṃvidaṃ karoti, tarhi kāryaṃ bhaviṣyati.}
\end{quote}

\begin{quote}
If you come to an agreement with him, the work will get done.
\end{quote}

A feminine noun.

\begin{cousinweb}
\textbf{\sk{संविद्}} (\emph{saṃvid}) is built on \emph{sam}, \textquotedblleft{}together\textquotedblright{}, and the root √vid, \textquotedblleft{}to know\textquotedblright{}: a knowing together; MW gives \emph{saṃvid} as consciousness, and as a mutual understanding, an agreement (MW).
\end{cousinweb}

\begin{grammarlens}[title={What you've built}]
One more word for saying what happens if.
\end{grammarlens}

\subsection*{Your turn}
\begin{itemize}
\raggedright
  \item \emph{Say it:} \emph{saṃvid}
  \item \emph{Say it:} \emph{saṃvid}, once more
  \item \emph{Say it:} \emph{yadi bhavān tena saha saṃvidaṃ karoti, tarhi kāryaṃ bhaviṣyati.} — If you come to an agreement with him, the work will get done.
  \item \emph{Recall:} say the Sanskrit for trouble, then the Sanskrit for a go-between, then say \emph{saṃvid} again
\end{itemize}

\begin{tcolorbox}[breakable,colback=teal!4,colframe=teal!35!black,title={Before you move on}]
What does \sk{संविद्} mean? (\textquotedblleft{}An agreement\textquotedblright{}.) Say it once more.
\end{tcolorbox}
\section[\texorpdfstring{\sk{अंगीकरोति} (aṅgīkaroti)}{अंगीकरोति (aṅgīkaroti)}]{\sk{अंगीकरोति} (aṅgīkaroti) — agrees to}

\label{lesson:SA-C534-angikaroti}



\begin{quote}
Before the new one: say the Sanskrit for a go-between, then the Sanskrit for an agreement.
\end{quote}

\subsection*{You'll want to know: \sk{अंगीकरोति}}
\textbf{\sk{अंगीकरोति}} — \emph{aṅgīkaroti} — \textquotedblleft{}agrees to\textquotedblright{}.

\begin{quote}
\sk{यदि} \sk{सः} \sk{मम} \sk{योजनाम्} \sk{अंगीकरोति}, \sk{तर्हि} \sk{वयं} \sk{श्वः} \sk{आरभामहे।}
\end{quote}

\begin{quote}
\emph{yadi saḥ mama yojanām aṅgīkaroti, tarhi vayaṃ śvaḥ ārabhāmahe.}
\end{quote}

\begin{quote}
If he agrees to my plan, we will start tomorrow.
\end{quote}

Past (he …): \sk{सः} \sk{योजनाम्} \sk{अंगीकृतवान्}, \emph{saḥ yojanām aṅgīkṛtavān}.

\begin{cousinweb}
\textbf{\sk{अंगीकरोति}} (\emph{aṅgīkaroti}) is built from \emph{aṅga}, \textquotedblleft{}a limb, a part\textquotedblright{}, in the form \emph{aṅgī}, and the root √kṛ, \textquotedblleft{}to make\textquotedblright{}: to make a thing part of oneself; MW gives \emph{aṅgīkṛ} as to take the side of, to agree to (MW).
\end{cousinweb}

\begin{grammarlens}[title={What you've built}]
One more word for saying what happens if.
\end{grammarlens}

\subsection*{Your turn}
\begin{itemize}
\raggedright
  \item \emph{Say it:} \emph{aṅgīkaroti}
  \item \emph{Say it:} \emph{aṅgīkaroti}, once more
  \item \emph{Say it:} \emph{yadi saḥ mama yojanām aṅgīkaroti, tarhi vayaṃ śvaḥ ārabhāmahe.} — If he agrees to my plan, we will start tomorrow.
  \item \emph{Recall:} say the Sanskrit for a go-between, then the Sanskrit for an agreement, then say \emph{aṅgīkaroti} again
\end{itemize}

\begin{tcolorbox}[breakable,colback=teal!4,colframe=teal!35!black,title={Before you move on}]
What does \sk{अंगीकरोति} mean? (\textquotedblleft{}Agrees to\textquotedblright{}.) Say it once more.
\end{tcolorbox}
\section[\texorpdfstring{\sk{प्रत्याख्याति} (pratyākhyāti)}{प्रत्याख्याति (pratyākhyāti)}]{\sk{प्रत्याख्याति} (pratyākhyāti) — turns down}

\label{lesson:SA-C534-pratyakhyati}



\begin{quote}
Before the new one: say the Sanskrit for an agreement, then the Sanskrit for agrees to.
\end{quote}

\subsection*{You'll want to know: \sk{प्रत्याख्याति}}
\textbf{\sk{प्रत्याख्याति}} — \emph{pratyākhyāti} — \textquotedblleft{}turns down\textquotedblright{}.

\begin{quote}
\sk{यदि} \sk{मूल्यम्} \sk{उचितं} \sk{न} \sk{अस्ति}, \sk{तर्हि} \sk{अहं} \sk{तत्} \sk{प्रत्याख्यामि।}
\end{quote}

\begin{quote}
\emph{yadi mūlyam ucitaṃ na asti, tarhi ahaṃ tat pratyākhyāmi.}
\end{quote}

\begin{quote}
If the price is not fair, I turn it down.
\end{quote}

Past (he …): \sk{सः} \sk{निमन्त्रणं} \sk{प्रत्याख्यातवान्}, \emph{saḥ nimantraṇaṃ pratyākhyātavān}.

\begin{cousinweb}
\textbf{\sk{प्रत्याख्याति}} (\emph{pratyākhyāti}) is built on \emph{prati}, \textquotedblleft{}back\textquotedblright{}, \emph{ā} and the root √khyā, \textquotedblleft{}to tell\textquotedblright{}; MW gives \emph{pratyākhyā} as to refuse, to reject, to deny (MW).
\end{cousinweb}

\begin{grammarlens}[title={What you've built}]
One more word for saying what happens if.
\end{grammarlens}

\subsection*{Your turn}
\begin{itemize}
\raggedright
  \item \emph{Say it:} \emph{pratyākhyāti}
  \item \emph{Say it:} \emph{pratyākhyāti}, once more
  \item \emph{Say it:} \emph{yadi mūlyam ucitaṃ na asti, tarhi ahaṃ tat pratyākhyāmi.} — If the price is not fair, I turn it down.
  \item \emph{Recall:} say the Sanskrit for an agreement, then the Sanskrit for agrees to, then say \emph{pratyākhyāti} again
\end{itemize}

\begin{tcolorbox}[breakable,colback=teal!4,colframe=teal!35!black,title={Before you move on}]
What does \sk{प्रत्याख्याति} mean? (\textquotedblleft{}Turns down\textquotedblright{}.) Say it once more.
\end{tcolorbox}
\section[\texorpdfstring{\sk{सम्मतिः} (sammatiḥ)}{सम्मतिः (sammatiḥ)}]{\sk{सम्मतिः} (sammatiḥ) — consent; approval}

\label{lesson:SA-C534-sammatih}



\begin{quote}
Before the new one: say the Sanskrit for agrees to, then the Sanskrit for turns down.
\end{quote}

\subsection*{You'll want to know: \sk{सम्मतिः}}
\textbf{\sk{सम्मतिः}} — \emph{sammatiḥ} — \textquotedblleft{}consent; approval\textquotedblright{}.

\begin{quote}
\sk{यदि} \sk{पितुः} \sk{सम्मतिः} \sk{अस्ति}, \sk{तर्हि} \sk{अहं} \sk{गमिष्यामि।}
\end{quote}

\begin{quote}
\emph{yadi pituḥ sammatiḥ asti, tarhi ahaṃ gamiṣyāmi.}
\end{quote}

\begin{quote}
If my father agrees, I will go.
\end{quote}

A feminine noun.

\begin{cousinweb}
\textbf{\sk{सम्मतिः}} (\emph{sammatiḥ}) is built on \emph{sam}, \textquotedblleft{}together\textquotedblright{}, and \emph{mati}, \textquotedblleft{}thought\textquotedblright{}, from the root √man, \textquotedblleft{}to think\textquotedblright{}: a thinking alike; MW gives \emph{sammati} as agreement, approval (MW).
\end{cousinweb}

\begin{grammarlens}[title={What you've built}]
One more word for saying what happens if.
\end{grammarlens}

\subsection*{Your turn}
\begin{itemize}
\raggedright
  \item \emph{Say it:} \emph{sammatiḥ}
  \item \emph{Say it:} \emph{sammatiḥ}, once more
  \item \emph{Say it:} \emph{yadi pituḥ sammatiḥ asti, tarhi ahaṃ gamiṣyāmi.} — If my father agrees, I will go.
  \item \emph{Recall:} say the Sanskrit for agrees to, then the Sanskrit for turns down, then say \emph{sammatiḥ} again
\end{itemize}

\begin{tcolorbox}[breakable,colback=teal!4,colframe=teal!35!black,title={Before you move on}]
What does \sk{सम्मतिः} mean? (\textquotedblleft{}Consent; approval\textquotedblright{}.) Say it once more.
\end{tcolorbox}
